% !TEX root = ../main.tex
\chapter{Indledning}
\label{ch:indledning}

\section{Baggrund}

Denne rapport er udarbejdet i forbindelse med praktikforløb hos
Teknologisk Institut som del af uddannelsen til automationsteknolog
ved UCL Erhvervsakademi og Professionshøjskole.

\begin{wrapfigure}[14]{r}{0.30\textwidth}
    \centering
    \vspace{0.4em}
    \includegraphics[width=0.28\textwidth]{IMG20260613210156.jpg}
    \caption{Wattson.}
    \label{fig:wattson-oversigt}
\end{wrapfigure}

Teknologisk Institut viderefører et projekt, som Energinet
startede, om en fuldt 3D-printet humanoid robot ved navn Wattson.
Robotten består af 3D-printede mekaniske dele og standardelektronik
(servoer, mikrocontrollere og en Jetson-baseret styreenhed). Den
drives af en eksisterende ROS\,2-software-stak.

Wattson er oprindeligt tænkt til vedligeholds- og
inspektionsopgaver i højspændingsanlæg. En menneskeformet robot kan
færdes i den eksisterende infrastruktur, uden at anlæggene skal
bygges om.

Den oprindelige robot blev udviklet over ca.~ét år af to personer.
Både mekanisk design, elektrisk integration og softwarearkitektur
ligger i et offentligt git-repository. Hardwaren er fuldt
3D-printet, så det er i princippet muligt at fremstille en kopi ud
fra de eksisterende filer.

I praksis kræver det dog en betydelig integrationsindsats, før de
printede dele, elektronikken og softwaren udgør et fungerende
system. Opgaven i praktikken har været at gennemføre den
integration: fremstille de mekaniske dele, samle dem, integrere
elektronikken, idriftsætte softwaren og dokumentere processen, så
resultatet kan bruges som udgangspunkt for videre udvikling hos
Teknologisk Institut.

\clearpage
\section{Problemformulering}

Hvordan kan en humanoid robot, baseret på eksisterende 3D-design og
software, fremstilles, samles og integreres til en funktionel
enhed, hvor mekaniske og elektriske systemer korrekt understøtter
robottens bevægelser?

Projektet omfatter:

\begin{itemize}
    \item 3D-print og efterbehandling af robottens komponenter
    \item Mekanisk samling af robottens delsystemer
    \item Elektrisk opbygning og korrekt wiring af aktuatorer og
          styresystem
    \item Integration og afprøvning af eksisterende software
\end{itemize}

Der arbejdes med at sikre korrekt samspil mellem delkomponenterne
og med mulighed for tilpasning af udvalgte komponenter og styring,
så funktionalitet og stabilitet forbedres.

\section{Afgrænsning}

Projektet er afgrænset til en humanoid robotarm som et selvstændigt
delsystem. Armen er valgt, fordi den indeholder de centrale
mekaniske og styringsmæssige udfordringer: bevægelse i flere led,
belastning af komponenter og koordinering af flere aktuatorer.

Hele den humanoide robot færdiggøres derfor ikke. Der arbejdes
modulært, så det færdige delsystem kan indgå i det samlede
robotsystem senere.

Følgende dele af robotten behandles ikke i rapporten:

\begin{itemize}
    \item Den modsatte arm, torso, hofteparti, ben og hoved
    \item Teleoperation via VR-udstyr (Quest~3 og ZED-kamera), som
          det eksisterende softwaresystem understøtter
    \item Invers kinematik og bevægelsesplanlægning på systemniveau
\end{itemize}

\section{Målgruppe}

Rapporten henvender sig til teknikere og fagpersoner inden for
automation og robotteknologi, som arbejder med opbygning,
integration og idriftsættelse af robotsystemer. Formidlingen er
praktisk og teknisk.

\clearpage
\section{Metode}

Projektet er gennemført som et idriftsættelsesforløb. Et
eksisterende design er omsat til en fungerende fysisk enhed gennem
fremstilling, samling, elektrisk integration og
softwarekonfiguration.

Arbejdet er foregået iterativt. Når et problem opstod under samling
eller test, blev årsagen diagnosticeret, og en løsning blev
implementeret. De vigtigste af disse hændelser er medtaget i
rapporten, fordi de viser de praktiske konsekvenser af at bygge med
3D-printede komponenter og af at integrere et eksisterende, kun
delvist dokumenteret softwaresystem.

Ud over selve idriftsættelsen er der udviklet en række tilføjelser
til det oprindelige system. De er lavet for at gøre det praktisk at
arbejde med robotten: stabil portbinding for de tilsluttede
mikrocontrollere og en samlet betjenings-GUI til opstart og demo.
Dertil kommer en commissioning-pakke med værktøjer til at måle og
dokumentere armens opførsel. Tilføjelserne er beskrevet i de
relevante kapitler sammen med de problemer, der motiverede dem.

Test og afprøvning af software-tilføjelserne er foregået på den
oprindelige Wattson-robot, der har stået til rådighed under hele
forløbet. Det har gjort det muligt at idriftsætte og verificere
ændringerne på en fungerende platform, mens replicaen blev
fremstillet og samlet. Når replicaen er færdig, kan den samme
softwarekonfiguration overføres til den.

\section{Rapportens opbygning}

Rapporten er opbygget som følger:

\begin{itemize}
    \item \textbf{Kapitel~\ref{ch:teori}} -- baggrundsviden:
          seriebus-servoer, ROS\,2 som kommunikationsplatform og
          egenskaberne ved 3D-printede konstruktionsdele.
    \item \textbf{Kapitel~\ref{ch:analyse}} -- det eksisterende
          system og begrundelse for afgrænsningen til én arm.
    \item \textbf{Kapitel~\ref{ch:mekanik}} -- mekanisk
          fremstilling og samling.
    \item \textbf{Kapitel~\ref{ch:el}} -- elektrisk integration og
          wiring.
    \item \textbf{Kapitel~\ref{ch:idrift}} -- softwareopsætning,
          kalibrering, commissioning-værktøjer og test.
    \item \textbf{Kapitel~\ref{ch:diskussion}} -- diskussion af de
          centrale valg og fund fra idriftsættelsen.
    \item \textbf{Kapitel~\ref{ch:konklusion}} -- besvarelse af
          problemformuleringen.
    \item \textbf{Kapitel~\ref{ch:perspektivering}} -- det videre
          arbejde.
\end{itemize}

